\chapter{FGA representability of the Picard scheme}
\label{chap:Picard_FGAPicRepresentability}
% archon:covers AlgebraicJacobian/Picard/FGAPicRepresentability.lean

% STRATEGY NOTE (iter-174). This chapter is the A.2.c assembly sub-row of
% Route~A (positive-genus arm of nonempty_jacobianWitness). It wires the
% Quot-scheme engine (\texttt{chap:Picard_QuotScheme}, A.2.b) and the
% \'etale-sheafified relative Picard functor (\texttt{chap:Picard_RelPicFunctor},
% A.1.c) into the FGA representability statement for \(\Pic_{C/k}\), following
% Kleiman, ``The Picard scheme'' \S 4 (\cite{kleiman-picard}). It is a small
% \emph{assembly} chapter: every nontrivial input is a \(\) pointer to a
% sibling chapter, and the work here is to thread the inputs through
% Grothendieck's existence theorem (\cite[\S 4, Main Theorem
% th:main]{kleiman-picard}) and its specialisation to a complete variety over a
% field (\cite[\S 4, Cor.~cor:algsch]{kleiman-picard}). Identity-component
% material (\Picz) is A.3 and lives in a separate sibling chapter --- it is
% \emph{not} scoped into this chapter.

% NOTE (planner directive iter-174). This chapter must produce the
% \emph{representable Picard scheme} \(\Pic_{C/k}\) as a \(k\)-group scheme,
% locally of finite type, that represents the \'etale sheafification
% \(\Pic^\sharp_{(C/k)\et}\) defined in A.1.c. The three deliverables are
% (i) a line-bundle / Quot translation expressed via the Abel map
% \(\mathrm{Div}_{C/k} \to \Pic_{C/k}\) (Kleiman \S 3, \S 4) which factorises through the
% Hilbert scheme \(\Hilb_{C/k} = \Quot_{\mathcal{O}_C/C/k}\), hence through the
% Quot engine of A.2.b; (ii) the existence theorem itself, in the
% field-of-definition specialisation \(S = \Spec k\), \(X = C\) a smooth proper
% geometrically integral curve; (iii) the group-scheme refinement carried by
% the tensor product of invertible sheaves on \(C \times_k T\). The Lean
% encoding is delegated to a Lean file
% \texttt{AlgebraicJacobian/Picard/FGAPicRepresentability.lean} whose
% construction is described in the final section but not written in this
% chapter.

\section{Setup and motivation}
\label{sec:fga_pic_setup}

Throughout this chapter, fix a base field \(k\) and a smooth proper geometrically
integral curve \(C\) over \(k\) (the curve carrying the Jacobian). The relative
Picard presheaf \(\Pic^\sharp_{C/k} : (\Sch/k)^{op} \to \Set\) of
\texttt{chap:Picard_LineBundlePullback} has already been packaged --- as a
set-valued functor at A.1.b, and refined to its abelian-group-valued
\'etale-sheafified version \(\Pic^\sharp_{(C/k)\et}\) at
\texttt{chap:Picard_RelPicFunctor} (A.1.c). The Quot scheme machinery of
\texttt{chap:Picard_QuotScheme} (A.2.b) supplies the representing engine: for
\(C/k\) projective and flat with integral geometric fibers, the Hilbert scheme
\(\Hilb_{C/k}\) and its open subscheme \(\mathrm{Div}_{C/k}\) of relative effective
divisors are both representable by quasi-projective \(k\)-schemes. Grothendieck's
existence theorem for \(\Pic\) assembles these two ingredients via the Abel
map \(\mathrm{Div}_{C/k} \to \Pic^\sharp_{(C/k)\et}\) to produce a representing scheme,
the \emph{Picard scheme} \(\Pic_{C/k}\). The present chapter packages this
assembly: a line-bundle / Quot correspondence, the FGA representability
theorem, and the abelian-group-scheme refinement.

% SOURCE: [Kleiman], ``The Picard scheme'', \S 4 (FGA Explained Ch.~9 \S 9.4);
% arXiv:math/0504020, p.~26. Read from
% references/kleiman-picard-src/kleiman-picard.tex.

\section{Line bundles and the Quot / Hilbert scheme}
\label{sec:fga_pic_line_bundle_quot}

The bridge from line bundles on \(C \times_k T\) to the Quot scheme is the
Abel map: a relative effective divisor \(D \subset C \times_k T\) over \(T\)
(an \(S\)-flat closed subscheme whose ideal is invertible) determines a closed
subscheme of \(C \times_k T\) flat over \(T\), hence a \(T\)-point of the Hilbert
functor \(\Hilb_{C/k}(T)\), which is the special case \(\Quot_{\mathcal{O}_C/C/k}\)
of the Quot functor of \texttt{chap:Picard_QuotScheme}; and it determines an
invertible sheaf \(\mathcal{O}_{C \times_k T}(D)\) on the relative curve,
representing a class in \(\Pic^\sharp_{(C/k)\et}(T)\).


% SOURCE QUOTE PROOF: (from references/kleiman-picard-src/kleiman-picard.tex,
% L1843-L1866, proof of th:repDiv)
% "Set \(H:=\IHilb_{X/S}\), and let \(W\subset X\x H\) be the universal
% (closed) subscheme, and \(q\:W\to H\) the projection.  Let \(V\) denote the
% set of points \(w\in W\) at which \(W\) is an effective divisor. Plainly \(V\)
% is open in \(W\).  Set \(Z:=q(W-V)\).  Then \(Z\) is closed because \(q\) is
% proper.  Set \(U:=H-Z\).  Then \(U\) is open, and \(q^{-1}U\) is an effective
% divisor in \(X\x U\).  In fact, since \(q\) is flat, \(q^{-1}U\) is a relative
% effective divisor in \(X\x U/U\).
% [\dots]"
\begin{proof}
  
  By \texttt{chap:Picard_QuotScheme}, the Hilbert functor \(\Hilb_{C/k} =
  \Quot_{\mathcal{O}_C/C/k}\) is representable by a \(k\)-scheme; this uses
  \(C/k\) projective and flat (the smooth proper curve hypothesis is
  stronger). Let \(W \subseteq C \times_k \Hilb_{C/k}\) be the universal
  closed subscheme and \(q : W \to \Hilb_{C/k}\) the projection. The locus
  \(V \subseteq W\) where \(W\) is locally an effective Cartier divisor (i.e.\
  cut out by one regular element) is open. The complement \(W \setminus V\)
  is closed in \(W\), and the image \(Z := q(W \setminus V) \subseteq
  \Hilb_{C/k}\) is closed because \(q\) is proper. Hence \(\mathrm{Div}_{C/k} :=
  \Hilb_{C/k} \setminus Z\) is open in \(\Hilb_{C/k}\) and represents the
  subfunctor of relative effective divisors: on \(\mathrm{Div}_{C/k}\) the
  restriction of \(W\) is flat and cut out by one regular element fiberwise,
  so by the standard relative-divisor criterion (Kleiman~\S 3,
  Lem.~3.4) it is a relative effective divisor on \(C \times_k
  \mathrm{Div}_{C/k} / \mathrm{Div}_{C/k}\).

  The Abel map sends a relative effective divisor \(D \subseteq C \times_k T\)
  to the invertible sheaf \(\mathcal{O}_{C \times_k T}(D) := \mathcal{I}_D^{-1}\)
  (the dual of the ideal sheaf, which is invertible by the divisor
  hypothesis). The assignment is functorial in \(T\): pullback of an ideal
  sheaf along \(\mathrm{id}_C \times_k g : C \times_k T' \to C \times_k T\)
  takes the ideal of \(D\) to the ideal of \(D_{T'}\) (since \(D\) is flat over
  \(T\), the schematic inverse image \(D_{T'}\) is a relative effective
  divisor on \(C \times_k T'/T'\), again by Kleiman~\S 3, Lem.~3.4), hence
  takes \(\mathcal{O}_{C \times_k T}(D)\) to \(\mathcal{O}_{C \times_k T'}(D_{T'})\).
  Composing with the canonical projection
  \(\Pic(C \times_k T) \to \Pic^\sharp_{C/k}(T)\) of
  \texttt{thm:relative_pic_quotient_well_defined} and then with the
  \'etale sheafification of \texttt{chap:Picard_RelPicFunctor} yields the
  natural map \(A_{C/k}(T) : \mathrm{Div}_{C/k}(T) \to \Pic^\sharp_{(C/k)\et}(T)\).
\end{proof}

\section{The FGA representability theorem}
\label{sec:fga_pic_representability}

We now state and outline the proof of Grothendieck's existence theorem
for the Picard scheme, in the specialisation needed for the Jacobian arm
of Route~A: \(S = \Spec k\), \(X = C\) a smooth proper geometrically integral
curve. The general theorem (Kleiman \S 4, Thm.~4.8) covers any
\(X/S\) projective Zariski-locally over \(S\), flat, with integral geometric
fibers; we record the general statement and then specialise.


% SOURCE QUOTE PROOF: (from references/kleiman-picard-src/kleiman-picard.tex,
% L2168-L2366, proof of th:main; long and is broken into the four-step
% sketch below per the project's structural restatement convention. The
% verbatim quote of the opening reduction step is:)
% "By~\cite[(0, 4.5.5), p.~106]{EGAG}, it is a local matter on \(S\) to
% represent a Zariski sheaf on the category of \(S\)-schemes.  Moreover, it
% is also a local matter on \(S\) to prove that an \(S\)-scheme is separated
% and locally of finite type.  Hence, in order to prove (1), we may assume
% \(S\) is Noetherian and \(X/S\) is projective.
% [\dots]
% Plainly an \(S\)-scheme is separated if it is a disjoint union of
% separated open subschemes, or if it is an increasing union of
% separated open subschemes.  Hence (1) follows from (2)."
% (subsequent steps quoted inline in the project-notation proof below).
\begin{proof}
  
  We follow Kleiman~\S 4, Thm.~4.8, specialised to \(S = \Spec k\) and
  \(X = C\). The proof proceeds in four conceptual steps; the underlying
  algebraic engine is the Quot scheme of \texttt{chap:Picard_QuotScheme}
  and the \'etale-sheafification of \texttt{chap:Picard_RelPicFunctor},
  together with the line-bundle / Quot correspondence of
  \texttt{lem:line_bundle_quot_correspondence}.

  \emph{Step 1: stratify by Hilbert polynomial.}
  For each polynomial \(\phi \in \mathbb{Q}[n]\), let \(P^\phi \subseteq
  \Pic^\sharp_{(C/k)\et}\) be the \'etale subsheaf whose \(T\)-points are
  represented by invertible sheaves \(\mathcal{L}\) on \(C \times_k T\) with
  \(\chi(C_t, \mathcal{L}_t^{-1}(n)) = \phi(n)\) for all \(t \in T\) (the
  Hilbert polynomial relative to the chosen polarisation
  \(\mathcal{O}_C(1)\) on the projective curve \(C\)). The \(P^\phi\) form a
  disjoint open cover of \(\Pic^\sharp_{(C/k)\et}\) (openness via Kleiman's
  flat-base-change argument; cf.\ EGA III 7.9.11). It suffices to
  represent each \(P^\phi\) by an increasing union of open quasi-projective
  \(k\)-schemes.

  \emph{Step 2: bound twists by \(m\)-regularity.}
  Given \(\phi\), fix \(m\) large enough that the \(T\)-points of \(P^\phi\) are
  representable by \(\mathcal{L}\) on \(C \times_k T\) satisfying
  \(R^i \pi_{T*} \mathcal{L}(n) = 0\) for all \(i \ge 1\), \(n \ge m\). Twisting
  by \(\mathcal{O}_C(m)\) defines an automorphism of \(\Pic^\sharp_{(C/k)\et}\)
  that carries \(P^\phi\) to \(P^{\phi_0}_0\) for \(\phi_0(n) := \phi(m + n)\).
  Hence it suffices to represent each \(P^{\phi_0}_0\) by a quasi-projective
  \(k\)-scheme.

  \emph{Step 3: factor through the Abel map.}
  Consider the Abel map \(A_{C/k} : \mathrm{Div}_{C/k} \to \Pic^\sharp_{(C/k)\et}\)
  of \texttt{lem:line_bundle_quot_correspondence}. The source \(\mathrm{Div}_{C/k}\)
  is representable by an open subscheme of \(\Hilb_{C/k}\) (and hence by an
  open subscheme of the Quot scheme of \texttt{thm:quot_representable}), and
  is quasi-projective over \(k\). Pulling back along
  \(P^{\phi_0}_0 \hookrightarrow \Pic^\sharp_{(C/k)\et}\) yields an open
  subscheme \(Z := P^{\phi_0}_0 \times_{\Pic^\sharp_{(C/k)\et}} \mathrm{Div}_{C/k}\)
  of \(\mathrm{Div}_{C/k}\), which is again quasi-projective. The projection
  \(\alpha : Z \to P^{\phi_0}_0\) is a surjection of \'etale sheaves:
  given a \(T\)-point \(\lambda\) of \(P^{\phi_0}_0\) represented by an \'etale
  cover \(T' \to T\) and an invertible sheaf \(\mathcal{L}'\) on \(C \times_k
  T'\) with \(H^1(C_{t'}, \mathcal{L}'_{t'}) = 0\), the projective bundle
  \(\mathbb{P}(\mathcal{Q}) \to T'\) associated to
  \(\mathcal{Q} := \pi_{T'*}\mathcal{L}'\) is smooth over \(T'\); an \'etale
  cover \(T_1 \to T'\) admits a \(T'\)-map \(T_1 \to \mathbb{P}(\mathcal{Q})\)
  (EGA IV 17.16.3 (ii)) whose composition with the universal divisor map
  \(\mathbb{P}(\mathcal{Q}) \to Z\) produces the required \(T_1\)-point of \(Z\).
  Moreover, the first projection \(Z \times_{P^{\phi_0}_0} Z \to Z\) is
  smooth and proper.

  \emph{Step 4: descend by the smooth-proper quotient lemma.}
  Apply Kleiman~\S 4, Lem.~4.9: given a surjection \(\alpha : Z \to P\)
  of \'etale sheaves with \(Z\) quasi-projective, \(R := Z \times_P Z\)
  representable, and the first projection \(R \to Z\) smooth and proper,
  \(P\) is representable by a quasi-projective scheme. In our case
  \(P = P^{\phi_0}_0\), \(Z\) and \(R\) as above, so \(P^{\phi_0}_0\) is
  quasi-projective. Reassembling over \(\phi \in \mathbb{Q}[n]\) and
  over \(m\), the disjoint union of representable \(P^\phi\) pieces represents
  \(\Pic^\sharp_{(C/k)\et}\). The total scheme is separated (as a disjoint
  union of separated quasi-projectives) and locally of finite type. This
  is Kleiman~\S 4, Cor.~4.18.3 in the \(S = \Spec k\),
  \(X = C\) complete case.
\end{proof}

\subsection{The smooth-proper quotient lemma}
\label{subsec:fga_pic_qt}

The structural lemma underlying step 4 above is recorded separately
because it is consumed verbatim in the assembly proof: it is the
abstract content allowing the descent from \(Z = \mathrm{Div}^{\phi_0}_0\) to the
representing Picard piece.


% SOURCE QUOTE PROOF: (from references/kleiman-picard-src/kleiman-picard.tex,
% L2377-L2417, proof of lm:qt; reproduced verbatim because it is the
% structural heart of step 4.)
% "To lighten the notation, let \(Z\) and \(R\) denote the
% corresponding schemes as well.  Since the structure map \(Z\to S\) is
% quasi-projective, it is separated; whence, the first projection
% \(Z\x_SZ\to Z\) is too.  But the first projection \(R\to Z\) is proper, and
% factors naturally through a map \(h\:R\to Z\x_SZ\).  Hence \(h\) is proper.
% But \(h\) is a monomorphism; that is, \(h\) is injective on \(T\)-points.
% Therefore, \(h\) is a closed embedding by \cite[8.11.5]{EGAIV3}.
%
% Plainly, for each \(S\)-scheme \(T\), the subset $R(T)\subset
% Z(T)\x_{S(T)}Z(T)$ is the graph of an equivalence relation.  Also, the
% map of schemes \(R\to Z\) is flat and proper, and \(Z\) is a
% quasi-projective \(S\)-scheme.  It follows that there exist a
% quasi-projective \(S\)-scheme \(Q\) and a faithfully flat and projective map
% \(Z\to Q\) such that \(R=Z\x_QZ\).  (In other words, a flat and proper
% equivalence relation on a quasi-projective scheme is effective.)
% [\dots]
% Since \(Z\to Q\) is flat, it is smooth if (and only if) its fibers are
% smooth.  But these fibers are, up to extension of the ground field, the
% same as those of \(R\to Z\).  And \(R\to Z\) is smooth by hypothesis.  Thus
% \(Z\to Q\) is smooth.
%
% It remains to see that \(Q\) represents \(P\).  First, observe that \(Z\to Q\)
% is (represents) a surjection of \'etale sheaves.
% [\dots]
% Since \(Z\to Q\) is a surjection of \'etale sheaves, \(Q\) is, in the
% category of \'etale sheaves, the coequalizer of the pair of maps
% \(R\rightrightarrows Z\) by Exercise~\texttt{ex:epi} below, which is an
% elementary exercise in general Grothendieck topology.  By the same
% exercise, \(P\) too is the coequalizer of this pair of maps.  But, in any
% category, the coequalizer is unique up to unique isomorphism.  Thus \(Q\)
% represents \(P\)."
\begin{proof}

\leanok
  By (ii) the structure map \(Z \to \Spec k\) is quasi-projective, hence
  separated, so the first projection \(Z \times_k Z \to Z\) is separated.
  By (iv) the first projection \(R \to Z\) is proper; it factors through
  \(h : R \to Z \times_k Z\) over the diagonal embedding (the second
  component is the second projection of \(R \to Z \times_P Z\)). The map
  \(h\) is therefore proper, and a monomorphism (because \(R\) is the graph
  of an equivalence relation on \(Z\), hence injective on \(T\)-points).
  By EGA IV 8.11.5, a proper monomorphism is a closed immersion: \(h\) is a
  closed embedding.

  Hence \(R \hookrightarrow Z \times_k Z\) is a flat (by (iv)) and proper
  equivalence relation on the quasi-projective \(k\)-scheme \(Z\). By the
  theory of flat-and-proper effective equivalence relations on
  quasi-projective schemes (Altman--Kleiman; the construction places
  the graph \(R \subseteq Z \times_k Z\) inside the universal subscheme of
  the Hilbert scheme \(\Hilb_{Z/k}\) and descends to a closed subscheme
  \(Q \subseteq \Hilb_{Z/k}\) representing the quotient), there exist a
  quasi-projective \(k\)-scheme \(Q\) and a faithfully flat, projective map
  \(Z \to Q\) with \(R = Z \times_Q Z\).

  The map \(Z \to Q\) is flat and proper; its fibers, up to field
  extension, coincide with those of \(R \to Z\), which are smooth by
  hypothesis. Hence \(Z \to Q\) is smooth.

  Finally, \(Z \to Q\) is a surjection of \'etale sheaves (an \'etale
  cover of any \(T\)-point of \(Q\) pulls back to a \(T'\)-point of \(Z\), by
  EGA IV 17.16.3 (ii) applied to \(Z \times_Q T \to T\)). In the category
  of \'etale sheaves, both \(Q\) and \(P\) are coequalisers of the pair
  \(R \rightrightarrows Z\); the coequaliser is unique up to unique
  isomorphism, so \(Q\) represents \(P\). Smoothness of \(\alpha : Z \to P\)
  is the smoothness of \(Z \to Q\) just shown.
\end{proof}

\section{Group-scheme structure}
\label{sec:fga_pic_group_structure}

The Picard scheme \(\Pic_{C/k}\) inherits an abelian group structure from
the tensor product of invertible sheaves on \(C \times_k T\). The relative
Picard presheaf \(T \mapsto \Pic(C \times_k T)/H_T\) is already an abelian
group at every \(T\) (the quotient of an abelian group by a subgroup),
functorially in \(T\); this lifts to the \'etale sheafification, and then
to a group-scheme structure on the representing scheme by Yoneda.


% SOURCE QUOTE PROOF: not applicable --- Kleiman \S 2--4 do not isolate a
% separate proof block for the group-scheme structure: it is implicit in
% the fact that \(\Pic_{(X/S)\et}\) takes values in abelian groups (df:Pfs)
% and that representability transfers algebraic operations from set-valued
% functors to scheme morphisms by Yoneda. The structural restatement below
% records the standard transfer.
\begin{proof}
  
  The relative Picard functor \(\Pic^\sharp_{C/k}\) of
  \texttt{thm:relative_pic_quotient_well_defined} is already abelian-group
  valued, with sum \([\mathcal{L}] + [\mathcal{M}] = [\mathcal{L} \otimes
  \mathcal{M}]\), inverse \(-[\mathcal{L}] = [\mathcal{L}^{-1}]\), and unit
  \(0 = [\mathcal{O}_{C \times_k T}]\); these operations descend through
  the quotient by \(H_T = \pi_T^*\Pic(T)\) because \(\pi_T^*\) is itself a
  group homomorphism. The \'etale sheafification step of
  \texttt{chap:Picard_RelPicFunctor} preserves the abelian-group structure
  on values, because sheafification of an abelian-presheaf is again an
  abelian-sheaf (the sheafification functor commutes with finite products
  in the value category).

  By \texttt{thm:fga_pic_representability}, \(\Pic^\sharp_{(C/k)\et}\) is
  represented by the scheme \(\Pic_{C/k}\). The Yoneda embedding is fully
  faithful and preserves limits, so the natural transformations
  $\Pic^\sharp_{(C/k)\et} \times \Pic^\sharp_{(C/k)\et} \to
  \Pic^\sharp_{(C/k)\et}$ (multiplication),
  \(\Pic^\sharp_{(C/k)\et} \to \Pic^\sharp_{(C/k)\et}\) (inverse), and the
  unit map \(h_{\Spec k} \to \Pic^\sharp_{(C/k)\et}\) pick out unique
  morphisms of representing schemes \(m\), \(i\), \(e\) as displayed. The
  abelian-group axioms (associativity, commutativity, identity,
  inverse) hold at the level of \(T\)-points functorially in \(T\), hence
  hold on the representing scheme by Yoneda. The local-finite-type
  property is part of \texttt{thm:fga_pic_representability}.
\end{proof}

\section{The Picard scheme: the named declaration}
\label{sec:fga_pic_named_definition}

The two preceding theorems are existence statements. We record the
named declaration corresponding to the Lean target.

\begin{definition}\leanok
[The Picard scheme]
  \label{def:pic_scheme}
  \lean{AlgebraicGeometry.Scheme.PicScheme}
  % NOTE iter-283: corrected this \uses to match the actual Lean definitional
  % order and break two dependency cycles that crashed `leanblueprint web`
  % (plastexdepgraph infinite recursion). In Lean (FGAPicRepresentability.lean
  % L223) `PicScheme := (HasPicScheme.has_pic_scheme C).choose`, i.e. it is
  % extracted from the `HasPicScheme` carrier (def:has_pic_scheme) ONLY -- NOT
  % from `representable` (thm:fga_pic_representability) and NOT from the group
  % structure (thm:pic_is_group_scheme). The two removed edges were both
  % spurious back-edges:
  %   (a) thm:pic_is_group_scheme: the abelian-group-scheme structure is built
  %       ON TOP of PicScheme (def:pic_scheme_group_object -> thm:pic_is_group_scheme,
  %       both \uses def:pic_scheme), so listing it here made a 2-cycle.
  %   (b) thm:fga_pic_representability: `representable` is extracted FROM the
  %       PicSharpRepresentable carrier whose statement references PicScheme
  %       (def:pic_sharp_representable -> def:pic_scheme, thm -> def:pic_sharp_representable),
  %       so listing it here made a 3-cycle. PicScheme depends only on HasPicScheme.
  % The explanatory \texttt{thm:fga_pic_representability}/\texttt{thm:pic_is_group_scheme}
  % in the prose are cross-references, not dependency edges, and are retained.
  \uses{def:has_pic_scheme, def:inst_has_pic_scheme}

  % SOURCE: [Kleiman], ``The Picard scheme'', \S 4, Def.~df:Psch;
  % arXiv:math/0504020, p.~26 (read from
  % references/kleiman-picard-src/kleiman-picard.tex, L2057-L2062).
  % SOURCE QUOTE:
  % "\begin{dfn}\label{df:Psch}
  %    If any of the four relative Picard functors of Definition~\texttt{df:Pfs}
  %  is representable, then the representing scheme is called the {\it Picard
  %  scheme\/} and denoted by \(\IPic_{X/S}\).  Moreover, we say simply that the
  %  Picard scheme \(\IPic_{X/S}\) exists."
  \textit{Source: [Kleiman], ``The Picard scheme'', \S 4, Def.~4.1.}
  Let \(C/k\) be a smooth proper geometrically integral curve over a field
  \(k\). The \emph{Picard scheme} \(\Pic_{C/k}\) is the \(k\)-scheme representing
  \(\Pic^\sharp_{(C/k)\et}\), supplied by
  \texttt{thm:fga_pic_representability}, equipped with the abelian-group-scheme
  structure of \texttt{thm:pic_is_group_scheme}. The scheme is separated,
  locally of finite type over \(k\), and a disjoint union of open
  quasi-projective \(k\)-subschemes.
\end{definition}

\section{Lean encoding}
\label{sec:fga_pic_lean_encoding}

The Lean target file is the new
\texttt{AlgebraicJacobian/Picard/FGAPicRepresentability.lean}, sibling to
\texttt{AlgebraicJacobian/Picard/RelativeSpec.lean}
(\texttt{chap:Picard_RelativeSpec}) and
\texttt{AlgebraicJacobian/Picard/LineBundlePullback.lean}
(\texttt{chap:Picard_LineBundlePullback}). The blueprint declarations of
this chapter correspond to Lean targets as follows.

\begin{itemize}
  \item \texttt{AlgebraicGeometry.Scheme.PicScheme}
    (\cref{def:pic_scheme}) is a structural constant naming the scheme
    representing \(\Pic^\sharp_{(C/k)\et}\); it is the scheme witness
    extracted from the \texttt{HasPicScheme} carrier
    (\cref{def:has_pic_scheme}) via \texttt{(HasPicScheme.has\_pic\_scheme
    C).choose}. The representability identification \texttt{representable}
    and the group-scheme structure are then layered on top of this bare
    scheme (so \cref{def:pic_scheme} depends only on \cref{def:has_pic_scheme},
    not on \texttt{thm:fga_pic_representability} or
    \texttt{thm:pic_is_group_scheme}).
  \item \texttt{AlgebraicGeometry.Scheme.PicScheme.abelMap}
    (\texttt{lem:line_bundle_quot_correspondence}) is the morphism of
    \(k\)-schemes \(\mathrm{Div}_{C/k} \to \Pic^\sharp_{(C/k)\et}\) sending a
    relative effective divisor to its associated invertible sheaf.
    The construction uses the Hilbert / Quot scheme produced by
    \texttt{chap:Picard_QuotScheme} (named
    \texttt{AlgebraicGeometry.Scheme.QuotScheme} per the directive's
    sibling-chapter naming), restricts to the open subscheme of relative
    effective divisors via the Kleiman~\S 3 effective-divisor open
    criterion, and composes with the quotient map to
    \(\Pic^\sharp_{(C/k)\et}\).
  \item \texttt{AlgebraicGeometry.Scheme.PicScheme.smoothProperQuotient}
    (\texttt{lem:smooth_proper_quotient}) is the Altman--Kleiman
    descent of a flat-and-proper equivalence relation \(R \subseteq Z
    \times_k Z\) on a quasi-projective \(k\)-scheme \(Z\) to a
    quasi-projective quotient \(Q\), packaged in the form that consumes
    a \(\Functor.\mathrm{RepresentableBy}\)-witness for \(Z\) and produces
    one for the coequalising sheaf.
  \item \texttt{AlgebraicGeometry.Scheme.PicScheme.representable}
    (\texttt{thm:fga_pic_representability}) is the assembly: from the
    \texttt{abelMap}, the \texttt{QuotScheme} representability witness,
    and the \texttt{rel\_pic\_etale\_sheafification} of
    \texttt{chap:Picard_RelPicFunctor}, build the
    \(\Functor.\mathrm{RepresentableBy}\)-witness for
    \(\Pic^\sharp_{(C/k)\et}\) via the four steps of the proof above
    (Hilbert-polynomial stratification, \(m\)-regularity bound, Abel-map
    factorisation, smooth-proper quotient).
  \item \texttt{AlgebraicGeometry.Scheme.PicScheme.groupSchemeStructure}
    (\texttt{thm:pic_is_group_scheme}) is the
    \texttt{CommGroupObject} (or
    \texttt{AddCommGrp\_Object}) instance on \texttt{PicScheme} obtained
    by Yoneda transport of the abelian-group structure on
    \(\Pic^\sharp_{(C/k)\et}\).
\end{itemize}

No new core axiom is required: this chapter is an assembly chapter on
top of the QuotScheme engine, the relative-Picard \'etale sheafification,
and the line-bundle pullback machinery already in place. The single
large external input (the Altman--Kleiman theorem of effective
flat-and-proper equivalence relations on quasi-projective schemes) is
recorded as \texttt{lem:smooth_proper_quotient} and is itself a standard
result; if its Lean formalisation is not yet in Mathlib, the
\texttt{smoothProperQuotient} Lean declaration may be sorried and the
gap reported to the plan agent, as it is logically downstream of the
Hilbert-scheme construction of \texttt{chap:Picard_QuotScheme}.

\section{Typeclass scaffolding for the representability assembly}
\label{sec:fga_pic_typeclass_scaffolding}

The five public declarations above are each phrased so that their single
mathematical obligation is isolated in a \emph{carrier}: a \texttt{Prop}-valued
existence predicate together with a noncomputable extraction of the asserted
object. This section records, with a one-to-one Lean correspondence, each
carrier predicate, each extraction, and each existence witness. The carriers
are file-internal scaffolding for the public results; their substantive
mathematical content is the obligation analysed in
\cref{sec:fga_pic_sorry_closure_order}.

\subsection{Carriers for the sibling-chapter functors}
\label{subsec:fga_pic_functor_carriers}

\begin{definition}\leanok
[Existence carrier for the relative Picard functor]
  \label{def:has_pic_sharp}
  \lean{AlgebraicGeometry.Scheme.PicScheme.HasPicSharp}
  The predicate \(\mathrm{HasPicSharp}(C)\) asserts that the category of presheaves of
  types on \((\Sch/k)^{op}\) admits at least one object intended to be the
  \'etale-sheafified relative Picard functor \(\Pic^\sharp_{(C/k)\et}\) of
  \texttt{chap:Picard_RelPicFunctor}. It is the existence carrier that lets the
  public declarations of this chapter mention \(\Pic^\sharp_{(C/k)\et}\)
  before the A.1.c sibling chapter has committed to a specific construction.
\end{definition}

\begin{proof}
  Proved directly in Lean.
\end{proof}

\begin{definition}\leanok
[The relative Picard functor as a carrier extraction]
  \label{def:pic_sharp_carrier}
  \lean{AlgebraicGeometry.Scheme.PicScheme.picSharp}
  \uses{def:has_pic_sharp, def:inst_has_pic_sharp}
  The presheaf \(\Pic^\sharp_{(C/k)\et} : (\Sch/k)^{op} \to \Set\) is the object whose
  existence is asserted by \cref{def:has_pic_sharp}, extracted as a chosen
  witness. It is a file-internal stand-in for the functor pinned in
  \texttt{chap:Picard_RelPicFunctor}; once that chapter lands, the extraction
  collapses to a re-export of \(\Pic^\sharp_{(C/k)\et}\).
\end{definition}

\begin{proof}
  Proved directly in Lean.
\end{proof}

\begin{definition}\leanok
[Existence witness for the relative Picard functor]
  \label{def:inst_has_pic_sharp}
  \lean{AlgebraicGeometry.Scheme.PicScheme.instHasPicSharp}
  \uses{def:has_pic_sharp}
  The canonical instance recording that \cref{def:has_pic_sharp} holds for
  every smooth proper geometrically integral curve \(C/k\). Its substantive
  discharge --- exhibiting the actual \'etale sheafification rather than a
  tautological witness --- is analysed as Sorry~1 in
  \cref{subsec:sorry_has_pic_sharp}.
\end{definition}

\begin{proof}

\leanok
  Proved directly in Lean; the semantic discharge is the obligation of
  \cref{subsec:sorry_has_pic_sharp}.
\end{proof}


\subsection{Carriers for the Picard scheme and its structure}
\label{subsec:fga_pic_scheme_carriers}

\begin{definition}\leanok
[Existence carrier for the Picard scheme]
  \label{def:has_pic_scheme}
  \lean{AlgebraicGeometry.Scheme.HasPicScheme}
  \uses{def:pic_sharp_carrier}
  The predicate \(\mathrm{HasPicScheme}(C)\) asserts that there exists a \(k\)-scheme
  \(X\) carrying a representability witness for \(\Pic^\sharp_{(C/k)\et}\); in
  other words, that the Picard scheme \(\Pic_{C/k}\) exists. It is the
  existence carrier from which \cref{def:pic_scheme} extracts the representing
  object.
\end{definition}

\begin{proof}
  Proved directly in Lean.
\end{proof}

\begin{definition}\leanok
[Existence witness for the Picard scheme]
  \label{def:inst_has_pic_scheme}
  \lean{AlgebraicGeometry.Scheme.instHasPicScheme}
  \uses{def:has_pic_scheme}
  The canonical instance recording that \cref{def:has_pic_scheme} holds for
  every smooth proper geometrically integral curve \(C/k\). Its substantive
  discharge --- the four-step FGA construction of
  \texttt{thm:fga_pic_representability}, gated on the Riemann--Roch substrate ---
  is analysed as Sorry~5 in \cref{subsec:sorry_has_pic_scheme}.
\end{definition}

\begin{proof}

\leanok
  Proved directly in Lean; the semantic discharge is the obligation of
  \cref{subsec:sorry_has_pic_scheme}.
\end{proof}


\section{Sorry-by-sorry closure order}
\label{sec:fga_pic_sorry_closure_order}

% iter-200 plan agent: stale-prose findings from iter-199
% lean-vs-blueprint-checker fga-iter199 addressed below (intro paragraph,
% Sorry 4 subsection location, closure-order summary Sorry 4 motivation).
The Lean target file
\texttt{AlgebraicJacobian/Picard/FGAPicRepresentability.lean} currently
carries seven open sorries, all of them \texttt{⟨sorry⟩} instance
bodies for the seven carrier typeclasses
\texttt{HasPicSharp}, \texttt{HasDivFunctor}, \texttt{HasPicScheme},
\texttt{HasAbelMap}, \texttt{HasSmoothProperQuotient},
\texttt{PicSharpRepresentable}, \texttt{PicSchemeGroupObject}.
After the iter-199 carrier-soundness refactor the file is
structurally homogeneous: every sorry is isolated in a single
\texttt{⟨sorry⟩} instance constructor of a \texttt{Prop}-valued
typeclass, so every theorem body extracts cleanly from the
appropriate field of the corresponding typeclass and is itself
axiom-clean. The carrier-soundness refactor of the
preceding section isolates each typeclass instance as the unique
sorry-carrying site for its associated extracted definition, so a
consumer that abstracts over the typeclass is kernel-clean. This
section is the explicit closure-order plan: for each sorry we record
(i) the Lean location and identifier, (ii) the mathematical content
to be proved, (iii) the Mathlib prerequisites at the project's pinned
revision, (iv) the Route~C (Riemann--Roch) substrate dependency, and
(v) a closure-order rank, where rank~1 is an independent carrier that
can be discharged from existing Mathlib alone, rank~2 depends on a
prior carrier or a Mathlib gap that is not Riemann--Roch, and rank~3
depends on a Route~C substrate re-engagement.

\medskip

The ranks are not a serialisation: rank-1 carriers may be closed in
parallel; rank-2 and rank-3 carriers depend either on other carriers
in the same chapter, on sibling A.1.c / A.2.b chapters
(\texttt{chap:Picard_RelPicFunctor},
\texttt{chap:Picard_QuotScheme}), or on the Route~C arc of the project
strategy (the Riemann--Roch development for smooth proper curves over
a field). The strategic point of the ranks is to identify which
sorries can be attempted first while the Route~C arc remains
substrate-blocked.

\subsection{Sorry 1 --- \texttt{instHasPicSharp} (L149)}
\label{subsec:sorry_has_pic_sharp}

\paragraph{Location.} Lean file
\texttt{AlgebraicJacobian/Picard/FGAPicRepresentability.lean},
line~149; instance declaration
\texttt{AlgebraicGeometry.Scheme.PicScheme.instHasPicSharp}.

\paragraph{Mathematical content.} The carrier asserts that the
\'etale-sheafified relative Picard functor $\Pic^\sharp_{(C/k)\et}$ of
\texttt{chap:Picard_RelPicFunctor} exists as a presheaf of types on
$(\Sch/k)^{op}$. The intended witness is the functor of \texttt{def:rel_pic_sharp}
specialised to $X = C$ and post-composed with \'etale sheafification:
the assignment $T \mapsto \Pic(C \times_k T) / \pi_T^* \Pic(T)$
followed by the associated sheaf construction in the \'etale topology.

\paragraph{Mathlib prerequisites.} The typeclass body unfolds to
\texttt{Nonempty ((Over (Spec (.of k)))${}^{op}$ \,$\rightsquigarrow$\, Type u)}.
\begin{itemize}
  \item \texttt{Functor}, \texttt{Nonempty}, \texttt{Over},
    \texttt{Spec}: \emph{exists} in Mathlib b80f227.
  \item The \'etale-sheafified relative Picard functor as a specific
    Mathlib object: \emph{missing}. Sheafification of presheaves on a
    site exists (\texttt{CategoryTheory.GrothendieckTopology.Plus},
    \texttt{Sheafify}), and the \'etale Grothendieck topology on
    \texttt{Scheme} is partially available, but the assembled
    \texttt{PicSharp} functor used downstream by
    \texttt{chap:Picard_RelPicFunctor} is a project-internal target,
    not a Mathlib name.
  \item Bridge from \texttt{Pic(C ${}\times_k$ T) / $\pi_T^* \Pic(T)$}
    to a presheaf of types: \emph{needs-bridge}. Mathlib has
    \texttt{Module.Picard} / \texttt{AlgebraicGeometry.Pic} (the
    invertible-sheaf functor) but the relative quotient functor on
    $(\Sch/k)^{op}$ that the chapter cites is the project's bespoke
    construction in \texttt{chap:Picard_LineBundlePullback}.
\end{itemize}

\paragraph{Route~C substrate dependency.} None at the level of the
typeclass statement: \texttt{Nonempty} of a non-empty type is
discharged by exhibiting any presheaf of types on
$(\Sch/k)^{op}$, e.g.\ the constant presheaf at \texttt{Unit}. The
semantic obligation --- that the chosen presheaf actually equals the
\'etale sheafification of \texttt{chap:Picard_RelPicFunctor} --- is not
captured by the \texttt{Prop}-valued typeclass. Closing the sorry to a
witness with the intended semantic identity requires the A.1.c sibling
chapter to land. No Riemann--Roch substrate is required.

\paragraph{Closure recipe (iter-199+).} After
\texttt{chap:Picard_RelPicFunctor} lands the pin
\texttt{AlgebraicGeometry.Scheme.PicSharp}, the present typeclass
instance reduces to \texttt{⟨⟨AlgebraicGeometry.Scheme.PicSharp C⟩⟩}
and the file-internal extraction \texttt{picSharp} collapses to a
one-line re-export. Until then a tautological witness (constant
presheaf) closes the typeclass without semantic content, which is
acceptable for the carrier-soundness probe but blocks downstream
mathematical use of \texttt{picSharp}.

\paragraph{Closure-order rank.} \textbf{Rank 2} (depends on the A.1.c
sibling chapter; no Riemann--Roch substrate).

\subsection{Sorry 2 --- \texttt{instHasDivFunctor} (L176)}
\label{subsec:sorry_has_div_functor}

\paragraph{Location.} Lean file
\texttt{AlgebraicJacobian/Picard/FGAPicRepresentability.lean},
line~176; instance
\texttt{AlgebraicGeometry.Scheme.PicScheme.instHasDivFunctor}.

\paragraph{Mathematical content.} The carrier asserts that the
relative-effective-divisor functor $\mathrm{Div}_{C/k}$ exists as a
presheaf of types on $(\Sch/k)^{op}$. The intended witness sends a
$k$-scheme $T$ to the set of relative effective Cartier divisors on
$C \times_k T$ flat over $T$; by \texttt{lem:line_bundle_quot_correspondence}
this functor is representable by an open subscheme of the Hilbert
scheme $\Hilb_{C/k} = \Quot_{\mathcal{O}_C / C / k}$.

\paragraph{Mathlib prerequisites.}
\begin{itemize}
  \item \texttt{Functor}, \texttt{Nonempty}, \texttt{Over},
    \texttt{Spec}: \emph{exists}.
  \item Relative effective Cartier divisor on a relative scheme:
    \emph{partial}. Mathlib has \texttt{Cartier} / \texttt{Divisor}
    machinery for divisors on a single scheme; the relative notion
    (flat over the base with invertible ideal sheaf) is partially
    spelled out in \texttt{AlgebraicGeometry.Divisor} but not as a
    presheaf-valued functor on the slice category.
  \item Hilbert scheme / Quot scheme open subfunctor of relative
    effective divisors (Kleiman \S3, Thm.~3.7):
    \emph{missing}. This is part of the A.2.b deliverable.
\end{itemize}

\paragraph{Route~C substrate dependency.} None at the typeclass
level; the same tautological discharge as Sorry~1 applies. The
semantic witness, however, requires the Quot/Hilbert engine of
\texttt{chap:Picard_QuotScheme}, which is currently substrate-blocked
on Route~C because the Quot construction itself uses Castelnuovo--
Mumford regularity bounds derived from cohomology and base change
(Nitsure~\S2; Mathlib lacks the relevant regularity machinery).

\paragraph{Closure recipe (iter-199+).} After
\texttt{chap:Picard_QuotScheme} lands the pin
\texttt{AlgebraicGeometry.Scheme.divFunctor} (the open subfunctor of
the Hilbert scheme cut out by the effective-divisor criterion), the
present typeclass instance reduces to
\texttt{⟨⟨AlgebraicGeometry.Scheme.divFunctor C⟩⟩}. Until then a
tautological closure leaves \texttt{divFunctor} semantically empty.

\paragraph{Closure-order rank.} \textbf{Rank 2} (depends on the
A.2.b sibling chapter; the Quot engine of that chapter is itself
Route~C-adjacent via Castelnuovo--Mumford regularity but does not
require Riemann--Roch for curves directly).

\subsection{Sorry 3 --- \texttt{instHasAbelMap} (L294)}
\label{subsec:sorry_has_abel_map}

\paragraph{Location.} Lean file
\texttt{AlgebraicJacobian/Picard/FGAPicRepresentability.lean},
line~294; instance
\texttt{AlgebraicGeometry.Scheme.PicScheme.instHasAbelMap}.

\paragraph{Mathematical content.} The carrier asserts existence of
the Abel map
$A_{C/k} : \mathrm{Div}_{C/k} \to \Pic^\sharp_{(C/k)\et}$ as a natural
transformation of presheaves on $(\Sch/k)^{op}$, sending a relative
effective Cartier divisor $D \subseteq C \times_k T$ to its associated
invertible sheaf $\mathcal{O}_{C \times_k T}(D) = \mathcal{I}_D^{-1}$.
This is the substantive content of
\texttt{lem:line_bundle_quot_correspondence}.

\paragraph{Mathlib prerequisites.}
\begin{itemize}
  \item Natural transformation \texttt{${\rightsquigarrow}$} of
    functors, \texttt{Nonempty}: \emph{exists}.
  \item Dual of an ideal sheaf $\mathcal{I}_D^{-1}$ as an invertible
    sheaf for a relative effective Cartier divisor: \emph{partial}.
    Mathlib has dual modules and dual sheaves but the specific link
    "ideal of a relative effective Cartier divisor is invertible and
    its dual is the associated line bundle" is a project-internal
    bridge.
  \item Functoriality of the assignment under pullback (Kleiman \S3,
    Lem.~3.4): \emph{needs-bridge}. Pullback of relative
    effective divisors along base-change is standard for flat-over-
    base divisors; Mathlib's coverage of this for the scheme-theoretic
    relative setting is incomplete.
\end{itemize}

\paragraph{Route~C substrate dependency.} None: the Abel map is a
construction at the level of line bundles and divisors, not a
consequence of Riemann--Roch. (Riemann--Roch tells us that the Abel
map for a fixed degree-$g$ pole-set is surjective onto $\Pic^g_{C/k}$;
this is a downstream consequence, not a prerequisite for defining the
map.)

\paragraph{Closure recipe (iter-199+).} The component-wise definition
of $A_{C/k}(T)$ is: given $D \in \mathrm{Div}_{C/k}(T)$, set
$A_{C/k}(T)(D) := [\mathcal{I}_D^{-1}]$ in
$\Pic^\sharp_{(C/k)\et}(T)$. Naturality follows from the pullback
formula for the dual ideal sheaf. The Lean encoding is a sequence of
three constructions: (a) extract the ideal sheaf of a relative
effective Cartier divisor, (b) dualise to get the line bundle,
(c) pass to the \'etale sheafified class. Each step is independent of
Riemann--Roch and depends only on Sorries~1 and~2 having committed to
a semantic witness for \texttt{picSharp} and \texttt{divFunctor}.

\paragraph{Closure-order rank.} \textbf{Rank 2} (depends on Sorries~1
and~2; no Riemann--Roch substrate).

\subsection{Sorry 4 --- \texttt{instHasSmoothProperQuotient} (L349)}
\label{subsec:sorry_smooth_proper_quotient}

\paragraph{Location.} Lean file
\texttt{AlgebraicJacobian/Picard/FGAPicRepresentability.lean},
line~349; instance constructor
\texttt{AlgebraicGeometry.Scheme.PicScheme.instHasSmoothProperQuotient}
of the iter-199 \texttt{Prop}-valued typeclass
\texttt{HasSmoothProperQuotient} (L338-341). The associated theorem
\texttt{AlgebraicGeometry.Scheme.PicScheme.smoothProperQuotient}
(L377) extracts \texttt{HasSmoothProperQuotient.is\_representable}
and is itself axiom-clean.

\paragraph{Mathematical content.} The smooth-proper quotient lemma
of \texttt{lem:smooth_proper_quotient}: given a surjection
$\alpha : Z \to P$ of \'etale sheaves on $(\Sch/k)^{op}$ with $Z$
quasi-projective, $R := Z \times_P Z$ representable, and the first
projection $R \to Z$ smooth and proper, the quotient $P$ is
representable by a quasi-projective $k$-scheme. The proof packages
the Altman--Kleiman effective-equivalence-relation theorem for
flat-and-proper relations on quasi-projective schemes.

\paragraph{Mathlib prerequisites.}
\begin{itemize}
  \item \texttt{Functor.IsRepresentable},
    \texttt{Functor.RepresentableBy}, \texttt{Limits.pullback}:
    \emph{exists}.
  \item \texttt{Smooth}, \texttt{IsProper} typeclasses on scheme
    morphisms: \emph{exists}.
  \item EGA IV~8.11.5 "proper monomorphism is a closed immersion":
    \emph{missing}. Mathlib has closed immersions, monomorphisms of
    schemes, and properness, but the EGA combination lemma is not yet
    isolated as a named result.
  \item Altman--Kleiman effective-equivalence-relation theorem
    (flat-and-proper equivalence relation on a quasi-projective
    scheme is effective, with quasi-projective quotient):
    \emph{missing}. This is the structural heart of the lemma and is
    not in Mathlib in any form.
  \item Coequaliser uniqueness in the category of \'etale sheaves:
    \emph{exists} (general categorical fact, uniqueness up to unique
    isomorphism of coequalisers in any category).
\end{itemize}

\paragraph{Route~C substrate dependency.} None: the lemma is a
structural fact about quasi-projective schemes and \'etale sheaves,
independent of curves and Riemann--Roch. It does, however, sit on the
Mathlib gap "effective equivalence relations of schemes", which is a
substantial development in its own right (corresponds to roughly
Altman--Kleiman, "Compactifying the Picard scheme", and the
flat-and-proper descent of equivalence relations from EGA IV~17).

\paragraph{Closure recipe (iter-199+).} Two routes are available:
\begin{enumerate}
  \item[(a)] Wait for the Mathlib gap "effective equivalence
    relations of quasi-projective schemes" to be filled and import
    the result. This is the cleanest long-term path but is outside
    the project's control timeline.
  \item[(b)] Formalise the Altman--Kleiman construction directly in
    the project as a sibling Lean file, gated on the EGA IV~8.11.5
    lemma (proper monomorphism is a closed immersion) which can be
    extracted from existing Mathlib infrastructure in a focused
    sub-task.
\end{enumerate}
Route~(b) is the recommended iter-199+ attempt: the EGA IV~8.11.5
extraction is bounded and the remaining structural construction
(graph of the equivalence relation inside the Hilbert scheme of $Z$,
descent to a closed subscheme of $\Hilb_{Z/k}$) is intrinsic.

\paragraph{Closure-order rank.} \textbf{Rank 2} (substantial Mathlib
gap; no Riemann--Roch substrate; recommended first attempt under the
carrier-soundness probe).

\subsection{Sorry 5 --- \texttt{instHasPicScheme} (L236)}
\label{subsec:sorry_has_pic_scheme}

\paragraph{Location.} Lean file
\texttt{AlgebraicJacobian/Picard/FGAPicRepresentability.lean},
line~236; instance
\texttt{AlgebraicGeometry.Scheme.PicScheme.instHasPicScheme}.

\paragraph{Mathematical content.} The carrier asserts that there
exists an object $X \in \mathrm{Over}(\Spec k)$ together with a
\texttt{RepresentableBy} witness for $\Pic^\sharp_{(C/k)\et}$. In
other words: the Picard scheme $\Pic_{C/k}$ exists. The semantic
content is the conclusion of Kleiman \S4, Thm.~4.8
specialised by Cor.~4.18.3 to $S = \Spec k$ and $X = C$
a smooth proper geometrically integral curve.

\paragraph{Mathlib prerequisites.}
\begin{itemize}
  \item \texttt{Functor.RepresentableBy}, \texttt{Nonempty},
    \texttt{Over}: \emph{exists}.
  \item \texttt{SmoothOfRelativeDimension 1},
    \texttt{IsProper}, \texttt{GeometricallyIntegral}: \emph{exists}
    as typeclasses, but the rich theory needed to actually produce a
    representing scheme is largely missing.
  \item The FGA representability machinery itself (the four-step
    construction of \texttt{thm:fga_pic_representability}, namely
    Hilbert polynomial stratification, $m$-regularity bound, Abel
    map factorisation, smooth-proper quotient): \emph{missing}.
  \item Castelnuovo--Mumford regularity / $m$-regularity for
    cohomologically twist-vanishing sheaves: \emph{missing}.
  \item Riemann--Roch for line bundles on a smooth proper curve over
    a field (used in the $m$-regularity step to bound twists by the
    degree): \emph{missing} (Route~C).
\end{itemize}

\paragraph{Route~C substrate dependency.} \textbf{Required}. The
$m$-regularity bound (Kleiman \S4, Step~2 of the proof of
Thm.~4.8) is the point at which the structural argument
meets Riemann--Roch: for a curve $C/k$ of genus $g$ and a line
bundle $\mathcal{L}$ of degree $d$, $H^1(C, \mathcal{L}(n)) = 0$ for
$n \geq m$ depends on $m \geq 2g - 1 - d + \deg \mathcal{O}_C(1)$, which
is the Serre vanishing consequence of Riemann--Roch. Specifically the
required result is the curve specialisation of
Hartshorne~III.5.2 (Serre vanishing) combined with IV.1.3
(Riemann--Roch for curves). Without Riemann--Roch, the $m$-regularity
step cannot be discharged and the FGA construction does not
terminate.

\paragraph{Closure recipe (iter-199+).} Conditional on the Route~C
re-engagement that the project STRATEGY.md flags as
substrate-blocked: once the Riemann--Roch arc lands, the $m$-regularity
step can be discharged via Hartshorne III.5.2 / IV.1.3, then the
remaining three steps (stratification, Abel-map factorisation,
smooth-proper quotient via Sorry~4) assemble into the existence
witness. Until then this sorry is the load-bearing blocker on the
Pic scheme arm of Route~A.

\paragraph{Closure-order rank.} \textbf{Rank 3} (requires the Route~C
substrate; load-bearing on the entire A.2.c representability claim).

\subsection{Sorry 6 --- \texttt{instPicSharpRepresentable} (L409)}
\label{subsec:sorry_pic_sharp_representable}

\paragraph{Location.} Lean file
\texttt{AlgebraicJacobian/Picard/FGAPicRepresentability.lean},
line~409; instance
\texttt{AlgebraicGeometry.Scheme.PicScheme.instPicSharpRepresentable}.

\paragraph{Mathematical content.} The carrier asserts that the
specific scheme \texttt{PicScheme C} (extracted from Sorry~5) is the
representing object of \texttt{picSharp C}: a witness in
\texttt{Nonempty ((picSharp C).RepresentableBy (PicScheme C))}. The
semantic content is the four-step assembly of
\texttt{thm:fga_pic_representability}: Hilbert polynomial
stratification of $\Pic^\sharp_{(C/k)\et}$ into open pieces
$P^\phi$, twist-bounding via $m$-regularity, factorisation through
the Abel map, descent via the smooth-proper quotient lemma.

\paragraph{Mathlib prerequisites.}
Identical to Sorry~5, with the additional dependency that the
representing object \texttt{PicScheme C} has been pinned by Sorry~5
(so the \texttt{RepresentableBy} witness has a destination object to
land in).
\begin{itemize}
  \item All Sorry~5 prerequisites apply (Castelnuovo--Mumford
    regularity, Riemann--Roch, Hilbert polynomial machinery):
    \emph{missing}.
  \item Hilbert polynomial stratification of an \'etale sheaf by
    Euler characteristic: \emph{missing}. Mathlib has Hilbert
    polynomials in commutative algebra but not the geometric stratification
    of a sheaf-of-points by a polynomial invariant.
  \item Abel map naturality and the factorisation
    $P^{\phi_0}_0 \times_{\Pic^\sharp} \mathrm{Div}_{C/k}$:
    \emph{needs-bridge} from Sorry~3.
  \item Smooth-proper quotient lemma: \emph{provided by Sorry~4}.
\end{itemize}

\paragraph{Route~C substrate dependency.} \textbf{Required}. Same
$m$-regularity bound as Sorry~5; the assembly is the proof of
representability itself, so it inherits every Riemann--Roch
dependency from Sorry~5.

\paragraph{Closure recipe (iter-199+).} Logically downstream of
Sorry~5 (which exhibits the representing object) and Sorry~4 (which
supplies the descent). After both prerequisite carriers are closed,
the four-step assembly is recorded in
\texttt{thm:fga_pic_representability} of this chapter and the Lean
target body extracts the \texttt{RepresentableBy} witness from the
existence statement.

\paragraph{Closure-order rank.} \textbf{Rank 3} (requires the
Route~C substrate via Sorry~5 and the structural Mathlib gap via
Sorry~4).

\subsection{Sorry 7 --- \texttt{instPicSchemeGroupObject} (L465)}
\label{subsec:sorry_pic_scheme_group_object}

\paragraph{Location.} Lean file
\texttt{AlgebraicJacobian/Picard/FGAPicRepresentability.lean},
line~465; instance
\texttt{AlgebraicGeometry.Scheme.PicScheme.instPicSchemeGroupObject}.

\paragraph{Mathematical content.} The carrier asserts existence of a
\texttt{GrpObj} (group-object) structure on \texttt{PicScheme C}.
The semantic content is the Yoneda transport of the abelian-group
operations $[\mathcal{L}] + [\mathcal{M}] = [\mathcal{L} \otimes \mathcal{M}]$,
$-[\mathcal{L}] = [\mathcal{L}^{-1}]$,
$0 = [\mathcal{O}_{C \times_k T}]$ from
$\Pic^\sharp_{(C/k)\et}(T)$ to morphisms of the representing scheme,
recorded in \texttt{thm:pic_is_group_scheme}.

\paragraph{Mathlib prerequisites.}
\begin{itemize}
  \item \texttt{GrpObj} (group-object structure in a category with
    finite products): \emph{exists}.
  \item \texttt{Functor.RepresentableBy} and Yoneda transport of
    natural transformations to scheme morphisms: \emph{exists}.
  \item Tensor product of invertible sheaves preserving the
    pullback / sheafification quotient: \emph{partial}; the
    project-internal lemma
    \texttt{thm:relative_pic_quotient_well_defined} of
    \texttt{chap:Picard_LineBundlePullback} carries the algebraic
    content.
  \item Sheafification of an abelian-group-valued presheaf commutes
    with the group structure: \emph{exists} (general fact about
    sheafification preserving finite products in the value category).
\end{itemize}

\paragraph{Route~C substrate dependency.} None directly; the
group-scheme structure transports from the value-level abelian-group
structure once the representing scheme exists. The dependency is
indirect via Sorry~6 (which requires the representing scheme to be
pinned and the \texttt{RepresentableBy} witness in place).

\paragraph{Closure recipe (iter-199+).} Conditional on Sorry~6
(the \texttt{RepresentableBy} witness for \texttt{picSharp C}): use
the Yoneda lemma to transport the abelian-group operations on
$\Pic^\sharp_{(C/k)\et}(T)$ to morphisms of \texttt{PicScheme C};
verify the group axioms at the level of $T$-points functorially in
$T$, hence on the representing scheme by Yoneda. The Lean
encoding does not require Riemann--Roch; it uses
\texttt{thm:pullback_natural} and
\texttt{thm:relative_pic_quotient_well_defined} from A.1.b together
with the \texttt{RepresentableBy} witness from Sorry~6.

\paragraph{Closure-order rank.} \textbf{Rank 3} (depends on Sorry~6,
which is Route~C-blocked, but is itself Route~C-independent
once Sorry~6 is in place).

\subsection{Closure-order summary}
\label{subsec:fga_pic_closure_order_summary}

The seven sorries partition into three ranks:
\begin{itemize}
  \item \textbf{Rank 1} (independent, closeable from existing
    Mathlib alone): none. Every sorry has either a sibling-chapter
    dependency, a Mathlib gap, or a Route~C substrate dependency.
  \item \textbf{Rank 2} (depends on sibling chapters or a
    non-Riemann--Roch Mathlib gap; closeable in iter-199+ under the
    carrier-soundness probe): Sorries~1, 2, 3, and~4.
    \begin{itemize}
      \item Sorries~1 and~2 are gated on the A.1.c
        (\texttt{chap:Picard_RelPicFunctor}) and A.2.b
        (\texttt{chap:Picard_QuotScheme}) sibling chapters
        committing to a specific \texttt{PicSharp} and
        \texttt{divFunctor} pin.
      \item Sorry~3 (Abel map) is gated on Sorries~1 and~2.
      \item Sorry~4 (smooth-proper quotient) is the substantive
        structural lemma; closeable independently of the other six
        by formalising the EGA IV~8.11.5 step and the Altman--Kleiman
        descent.
    \end{itemize}
  \item \textbf{Rank 3} (requires the Route~C substrate; blocked
    until the Riemann--Roch arc lands): Sorries~5, 6, and~7.
    \begin{itemize}
      \item Sorry~5 (Picard scheme existence) is the load-bearing
        Route~C-blocked carrier via the $m$-regularity / Serre
        vanishing dependency.
      \item Sorry~6 (representability witness) inherits the Route~C
        dependency of Sorry~5.
      \item Sorry~7 (group-scheme structure) is Route~C-independent
        in itself but depends on Sorry~6.
    \end{itemize}
\end{itemize}

\medskip

\textbf{Recommended first attempt under the carrier-soundness
probe.} Sorry~4 (the \texttt{instHasSmoothProperQuotient}
\texttt{⟨sorry⟩}) is the highest-value rank-2 target: it is a
structural lemma whose discharge does not depend on any sibling
chapter committing first, and after the iter-199 carrier-soundness
refactor it isolates the only mathematical obligation of
\texttt{lem:smooth_proper_quotient} into a single instance
constructor. The EGA IV~8.11.5 sub-step is bounded, and the
Altman--Kleiman descent has a known proof sketch in
\texttt{lem:smooth_proper_quotient} of this chapter.

\textbf{Recommended second tier under the carrier-soundness probe.}
Sorries~1, 2, 3 (the carrier instances for
\texttt{PicSharp}, \texttt{divFunctor}, \texttt{abelMap}) are
discharged in lockstep once the sibling chapters
\texttt{chap:Picard_RelPicFunctor} and \texttt{chap:Picard_QuotScheme}
commit to their pins. The dispatch order across the three is
1~$\to$~2~$\to$~3 (the Abel map needs both presheaves to exist).

\textbf{Held until Route~C re-engagement.} Sorries~5, 6, 7 form a
single dependency chain (5~$\to$~6~$\to$~7) all gated on the
Riemann--Roch arc. The Route~C re-engagement is the strategic
prerequisite for any progress on these three carriers; until it
lands, the rank-2 work suffices to keep the carrier-soundness probe
structurally complete and to allow downstream Route~A consumers to
proceed under typeclass quantification.

\section{Out of scope}
\label{sec:fga_pic_out_of_scope}

This chapter packages \emph{only} the representable Picard scheme
\(\Pic_{C/k}\) as an assembly of sibling-chapter inputs, plus the
group-scheme refinement. The following downstream constructions live in
separate sibling chapters and are explicitly out of scope here:
\begin{itemize}
  \item \textbf{Identity component \(\Pic^0_{C/k}\)} (A.3) --- the open and
    closed subgroup scheme of \(\Pic_{C/k}\) containing the unit and
    consisting of those classes lying in the connected component of
    \(\Pic_{C/k}\). Its dimension, smoothness, and abelian-variety
    structure are the content of the A.3 chapter
    \texttt{Picard\_Pic0AbelianVariety.tex} (planned), gated on this
    chapter.
  \item \textbf{Quot scheme representability} --- consumed by
    \texttt{thm:fga_pic_representability} via \texttt{thm:quot\_representable}
    (\texttt{chap:Picard_QuotScheme}), not re-proved here.
  \item \textbf{\'Etale sheafification of \(\Pic^\sharp\)} ---
    \texttt{chap:Picard_RelPicFunctor} (A.1.c) produces
    \(\Pic^\sharp_{(C/k)\et}\) from the set-valued presheaf of
    \texttt{chap:Picard_LineBundlePullback}; this chapter consumes the
    sheaf as a black box.
  \item \textbf{Universal Poincar\'e sheaf} (Kleiman~\S 4, Ex.~4.3)
    --- exists when \(\Pic_{C/k}\) represents \(\Pic_{C/k}\) (as opposed to
    its \'etale sheafification) and \(C/k\) has a section; out of scope
    here. The Albanese / Abel--Jacobi consumption in
    \texttt{chap:AbelJacobi} works with the \'etale-sheafified
    \(\Pic^\sharp_{(C/k)\et}\) representation and does not require the
    Poincar\'e sheaf.
  \item \textbf{Mumford's example} (Kleiman~\S 4, Eg.~4.14), counter-examples
    to representability for non-integral geometric fibers --- not
    relevant to a smooth proper geometrically integral curve.
\end{itemize}

\section{Internal-consistency check}
\label{sec:fga_pic_consistency_check}

The five declaration blocks form a connected \(\uses\) chain rooted in the
sibling chapters \texttt{chap:Picard_LineBundlePullback},
\texttt{chap:Picard_QuotScheme}, and \texttt{chap:Picard_RelPicFunctor}:
\begin{itemize}
  \item \texttt{lem:line_bundle_quot_correspondence} uses
    \texttt{def:line_bundle_on_product} (A.1.b) and
    \texttt{thm:pullback_natural} (A.1.b).
  \item \texttt{lem:smooth_proper_quotient} is a structural Archon-internal
    statement; no external \texttt{\textbackslash uses} pointers.
  \item \texttt{thm:fga_pic_representability} uses
    \cref{def:pic_scheme}, \texttt{lem:line_bundle_quot_correspondence},
    \texttt{thm:quot\_representable} (A.2.b),
    \texttt{def:rel\_pic\_etale\_sheafification} (A.1.c),
    \texttt{thm:pullback_natural} (A.1.b),
    \texttt{thm:relative_pic_quotient_well_defined} (A.1.b), and
    \texttt{lem:smooth_proper_quotient}.
  \item \texttt{thm:pic_is_group_scheme} uses
    \cref{def:pic_scheme}, \texttt{thm:fga_pic_representability},
    \texttt{def:rel\_pic\_etale\_sheafification} (A.1.c), and
    \texttt{thm:relative_pic_quotient_well_defined} (A.1.b).
  \item \cref{def:pic_scheme} uses
    \texttt{thm:fga_pic_representability} and
    \texttt{def:rel\_pic\_etale\_sheafification} (A.1.c).
\end{itemize}
The labels \texttt{thm:quot\_representable} and
\texttt{def:rel\_pic\_etale\_sheafification} are forward references to
the sibling chapters \texttt{Picard\_QuotScheme.tex} and
\texttt{Picard\_RelPicFunctor.tex} also being written this iteration; if
either lands with a different label, the plan agent should regenerate
the \texttt{\textbackslash uses} pointers across the three chapters in lockstep.
